% Options for packages loaded elsewhere
% Options for packages loaded elsewhere
\PassOptionsToPackage{unicode}{hyperref}
\PassOptionsToPackage{hyphens}{url}
\PassOptionsToPackage{dvipsnames,svgnames,x11names}{xcolor}
%
\documentclass[
  letterpaper,
  DIV=11,
  numbers=noendperiod]{scrartcl}
\usepackage{xcolor}
\usepackage{amsmath,amssymb}
\setcounter{secnumdepth}{-\maxdimen} % remove section numbering
\usepackage{iftex}
\ifPDFTeX
  \usepackage[T1]{fontenc}
  \usepackage[utf8]{inputenc}
  \usepackage{textcomp} % provide euro and other symbols
\else % if luatex or xetex
  \usepackage{unicode-math} % this also loads fontspec
  \defaultfontfeatures{Scale=MatchLowercase}
  \defaultfontfeatures[\rmfamily]{Ligatures=TeX,Scale=1}
\fi
\usepackage{lmodern}
\ifPDFTeX\else
  % xetex/luatex font selection
\fi
% Use upquote if available, for straight quotes in verbatim environments
\IfFileExists{upquote.sty}{\usepackage{upquote}}{}
\IfFileExists{microtype.sty}{% use microtype if available
  \usepackage[]{microtype}
  \UseMicrotypeSet[protrusion]{basicmath} % disable protrusion for tt fonts
}{}
\makeatletter
\@ifundefined{KOMAClassName}{% if non-KOMA class
  \IfFileExists{parskip.sty}{%
    \usepackage{parskip}
  }{% else
    \setlength{\parindent}{0pt}
    \setlength{\parskip}{6pt plus 2pt minus 1pt}}
}{% if KOMA class
  \KOMAoptions{parskip=half}}
\makeatother
% Make \paragraph and \subparagraph free-standing
\makeatletter
\ifx\paragraph\undefined\else
  \let\oldparagraph\paragraph
  \renewcommand{\paragraph}{
    \@ifstar
      \xxxParagraphStar
      \xxxParagraphNoStar
  }
  \newcommand{\xxxParagraphStar}[1]{\oldparagraph*{#1}\mbox{}}
  \newcommand{\xxxParagraphNoStar}[1]{\oldparagraph{#1}\mbox{}}
\fi
\ifx\subparagraph\undefined\else
  \let\oldsubparagraph\subparagraph
  \renewcommand{\subparagraph}{
    \@ifstar
      \xxxSubParagraphStar
      \xxxSubParagraphNoStar
  }
  \newcommand{\xxxSubParagraphStar}[1]{\oldsubparagraph*{#1}\mbox{}}
  \newcommand{\xxxSubParagraphNoStar}[1]{\oldsubparagraph{#1}\mbox{}}
\fi
\makeatother

\usepackage{color}
\usepackage{fancyvrb}
\newcommand{\VerbBar}{|}
\newcommand{\VERB}{\Verb[commandchars=\\\{\}]}
\DefineVerbatimEnvironment{Highlighting}{Verbatim}{commandchars=\\\{\}}
% Add ',fontsize=\small' for more characters per line
\usepackage{framed}
\definecolor{shadecolor}{RGB}{241,243,245}
\newenvironment{Shaded}{\begin{snugshade}}{\end{snugshade}}
\newcommand{\AlertTok}[1]{\textcolor[rgb]{0.68,0.00,0.00}{#1}}
\newcommand{\AnnotationTok}[1]{\textcolor[rgb]{0.37,0.37,0.37}{#1}}
\newcommand{\AttributeTok}[1]{\textcolor[rgb]{0.40,0.45,0.13}{#1}}
\newcommand{\BaseNTok}[1]{\textcolor[rgb]{0.68,0.00,0.00}{#1}}
\newcommand{\BuiltInTok}[1]{\textcolor[rgb]{0.00,0.23,0.31}{#1}}
\newcommand{\CharTok}[1]{\textcolor[rgb]{0.13,0.47,0.30}{#1}}
\newcommand{\CommentTok}[1]{\textcolor[rgb]{0.37,0.37,0.37}{#1}}
\newcommand{\CommentVarTok}[1]{\textcolor[rgb]{0.37,0.37,0.37}{\textit{#1}}}
\newcommand{\ConstantTok}[1]{\textcolor[rgb]{0.56,0.35,0.01}{#1}}
\newcommand{\ControlFlowTok}[1]{\textcolor[rgb]{0.00,0.23,0.31}{\textbf{#1}}}
\newcommand{\DataTypeTok}[1]{\textcolor[rgb]{0.68,0.00,0.00}{#1}}
\newcommand{\DecValTok}[1]{\textcolor[rgb]{0.68,0.00,0.00}{#1}}
\newcommand{\DocumentationTok}[1]{\textcolor[rgb]{0.37,0.37,0.37}{\textit{#1}}}
\newcommand{\ErrorTok}[1]{\textcolor[rgb]{0.68,0.00,0.00}{#1}}
\newcommand{\ExtensionTok}[1]{\textcolor[rgb]{0.00,0.23,0.31}{#1}}
\newcommand{\FloatTok}[1]{\textcolor[rgb]{0.68,0.00,0.00}{#1}}
\newcommand{\FunctionTok}[1]{\textcolor[rgb]{0.28,0.35,0.67}{#1}}
\newcommand{\ImportTok}[1]{\textcolor[rgb]{0.00,0.46,0.62}{#1}}
\newcommand{\InformationTok}[1]{\textcolor[rgb]{0.37,0.37,0.37}{#1}}
\newcommand{\KeywordTok}[1]{\textcolor[rgb]{0.00,0.23,0.31}{\textbf{#1}}}
\newcommand{\NormalTok}[1]{\textcolor[rgb]{0.00,0.23,0.31}{#1}}
\newcommand{\OperatorTok}[1]{\textcolor[rgb]{0.37,0.37,0.37}{#1}}
\newcommand{\OtherTok}[1]{\textcolor[rgb]{0.00,0.23,0.31}{#1}}
\newcommand{\PreprocessorTok}[1]{\textcolor[rgb]{0.68,0.00,0.00}{#1}}
\newcommand{\RegionMarkerTok}[1]{\textcolor[rgb]{0.00,0.23,0.31}{#1}}
\newcommand{\SpecialCharTok}[1]{\textcolor[rgb]{0.37,0.37,0.37}{#1}}
\newcommand{\SpecialStringTok}[1]{\textcolor[rgb]{0.13,0.47,0.30}{#1}}
\newcommand{\StringTok}[1]{\textcolor[rgb]{0.13,0.47,0.30}{#1}}
\newcommand{\VariableTok}[1]{\textcolor[rgb]{0.07,0.07,0.07}{#1}}
\newcommand{\VerbatimStringTok}[1]{\textcolor[rgb]{0.13,0.47,0.30}{#1}}
\newcommand{\WarningTok}[1]{\textcolor[rgb]{0.37,0.37,0.37}{\textit{#1}}}

\usepackage{longtable,booktabs,array}
\usepackage{calc} % for calculating minipage widths
% Correct order of tables after \paragraph or \subparagraph
\usepackage{etoolbox}
\makeatletter
\patchcmd\longtable{\par}{\if@noskipsec\mbox{}\fi\par}{}{}
\makeatother
% Allow footnotes in longtable head/foot
\IfFileExists{footnotehyper.sty}{\usepackage{footnotehyper}}{\usepackage{footnote}}
\makesavenoteenv{longtable}
\usepackage{graphicx}
\makeatletter
\newsavebox\pandoc@box
\newcommand*\pandocbounded[1]{% scales image to fit in text height/width
  \sbox\pandoc@box{#1}%
  \Gscale@div\@tempa{\textheight}{\dimexpr\ht\pandoc@box+\dp\pandoc@box\relax}%
  \Gscale@div\@tempb{\linewidth}{\wd\pandoc@box}%
  \ifdim\@tempb\p@<\@tempa\p@\let\@tempa\@tempb\fi% select the smaller of both
  \ifdim\@tempa\p@<\p@\scalebox{\@tempa}{\usebox\pandoc@box}%
  \else\usebox{\pandoc@box}%
  \fi%
}
% Set default figure placement to htbp
\def\fps@figure{htbp}
\makeatother

\ifLuaTeX
  \usepackage{luacolor}
  \usepackage[soul]{lua-ul}
\else
  \usepackage{soul}
\fi




\setlength{\emergencystretch}{3em} % prevent overfull lines

\providecommand{\tightlist}{%
  \setlength{\itemsep}{0pt}\setlength{\parskip}{0pt}}



 


\KOMAoption{captions}{tableheading}
\makeatletter
\@ifpackageloaded{caption}{}{\usepackage{caption}}
\AtBeginDocument{%
\ifdefined\contentsname
  \renewcommand*\contentsname{Table of contents}
\else
  \newcommand\contentsname{Table of contents}
\fi
\ifdefined\listfigurename
  \renewcommand*\listfigurename{List of Figures}
\else
  \newcommand\listfigurename{List of Figures}
\fi
\ifdefined\listtablename
  \renewcommand*\listtablename{List of Tables}
\else
  \newcommand\listtablename{List of Tables}
\fi
\ifdefined\figurename
  \renewcommand*\figurename{Figure}
\else
  \newcommand\figurename{Figure}
\fi
\ifdefined\tablename
  \renewcommand*\tablename{Table}
\else
  \newcommand\tablename{Table}
\fi
}
\@ifpackageloaded{float}{}{\usepackage{float}}
\floatstyle{ruled}
\@ifundefined{c@chapter}{\newfloat{codelisting}{h}{lop}}{\newfloat{codelisting}{h}{lop}[chapter]}
\floatname{codelisting}{Listing}
\newcommand*\listoflistings{\listof{codelisting}{List of Listings}}
\makeatother
\makeatletter
\makeatother
\makeatletter
\@ifpackageloaded{caption}{}{\usepackage{caption}}
\@ifpackageloaded{subcaption}{}{\usepackage{subcaption}}
\makeatother
\usepackage{bookmark}
\IfFileExists{xurl.sty}{\usepackage{xurl}}{} % add URL line breaks if available
\urlstyle{same}
\hypersetup{
  pdftitle={1.2: Intro to Coding in R},
  pdfauthor={Ellen Bledsoe},
  colorlinks=true,
  linkcolor={blue},
  filecolor={Maroon},
  citecolor={Blue},
  urlcolor={Blue},
  pdfcreator={LaTeX via pandoc}}


\title{1.2: Intro to Coding in R}
\author{Ellen Bledsoe}
\date{}
\begin{document}
\maketitle


\section{Introduction to Coding}\label{introduction-to-coding}

In this lesson, we will create objects, build vectors, subset data, use
functions, and work with data frames.

\subsection{Learning Outcomes}\label{learning-outcomes}

\begin{itemize}
\tightlist
\item
  Students will be able to define the following terms: object,
  assignment, vector, function, data frame
\item
  Students will be able to run code line-by-line and as code chunks from
  an Rmarkdown file.
\item
  Students will be able to write code assigning values to variables and
  use these variables to perform various operations.
\item
  Students will be able to recall and explain how functions operate, and
  the basic syntax around functions (arguments, auto-completion,
  parentheses).
\item
  Students will be able to differentiate different data classes in R.
\end{itemize}

\subsection{Assigning Objects}\label{assigning-objects}

An object is simply a name that stores a value in R so that we can reuse
it later.

Assignments are really key to almost everything we do in R. This is how
we create permanence in R. Anything can be saved to an object, and we do
this with the assignment operator, \texttt{\textless{}-}.

The short-cut for \texttt{\textless{}-} is \texttt{Alt\ +\ -} (or
\texttt{Option\ +\ -} on a Mac)

\begin{Shaded}
\begin{Highlighting}[]
\CommentTok{\# Assigning Objects}
\NormalTok{height }\OtherTok{\textless{}{-}} \FloatTok{47.5}           
\NormalTok{age }\OtherTok{\textless{}{-}} \DecValTok{122}

\CommentTok{\# We can do math with objects}
\NormalTok{height }\OtherTok{\textless{}{-}}\NormalTok{ height }\SpecialCharTok{*} \DecValTok{2}    \CommentTok{\# multiply}
\CommentTok{\# Notice how we are overwriting the original value of height. R replaces the old value with the new one.}
\NormalTok{age }\OtherTok{\textless{}{-}}\NormalTok{ age }\SpecialCharTok{{-}} \DecValTok{20}         \CommentTok{\# subtract}
\NormalTok{height\_index }\OtherTok{\textless{}{-}}\NormalTok{ height}\SpecialCharTok{/}\NormalTok{age  }\CommentTok{\# divide}
\NormalTok{height\_sq }\OtherTok{\textless{}{-}}\NormalTok{ height}\SpecialCharTok{\^{}}\DecValTok{2}       \CommentTok{\# raise to an exponent}

\CommentTok{\# This is simplistic and you\textquotesingle{}ll rarely do it in real{-}world scenarios. }
\end{Highlighting}
\end{Shaded}

\subsection{1-Dimensional Data: Vectors}\label{dimensional-data-vectors}

We can also assign more complex group of elements of the same type to a
particular object. This is called a \textbf{vector}, a basic data
structure in R.

All elements in a vector must be the same type (all numbers, all text,
etc.).

\begin{Shaded}
\begin{Highlighting}[]
\NormalTok{weight\_kg }\OtherTok{\textless{}{-}} \FunctionTok{c}\NormalTok{(}\DecValTok{3}\NormalTok{, }\DecValTok{2}\NormalTok{, }\DecValTok{4}\NormalTok{, }\DecValTok{9}\NormalTok{, }\DecValTok{7}\NormalTok{, }\DecValTok{3}\NormalTok{, }\DecValTok{6}\NormalTok{)}
\NormalTok{weight\_kg}
\end{Highlighting}
\end{Shaded}

\begin{verbatim}
[1] 3 2 4 9 7 3 6
\end{verbatim}

\begin{Shaded}
\begin{Highlighting}[]
\NormalTok{animals }\OtherTok{\textless{}{-}} \FunctionTok{c}\NormalTok{(}\StringTok{"cat"}\NormalTok{, }\StringTok{"rat"}\NormalTok{, }\StringTok{"bat"}\NormalTok{, }\StringTok{"rat"}\NormalTok{)}
\NormalTok{animals}
\end{Highlighting}
\end{Shaded}

\begin{verbatim}
[1] "cat" "rat" "bat" "rat"
\end{verbatim}

\begin{Shaded}
\begin{Highlighting}[]
\CommentTok{\# R does everything in vectors}
\end{Highlighting}
\end{Shaded}

\subsection{Data classes}\label{data-classes}

There are a few main types of data in R, and they behave differently. We
call these types of data ``classes.''

\begin{itemize}
\tightlist
\item
  numeric / double (numbers, decimals allowed)
\item
  integer (no decimals allowed)
\item
  character (letters or mixture)
\item
  logical (True or False; T or F)
\item
  factors (best used for data that need to be in a specific order;
  levels indicate the order)
\end{itemize}

\begin{Shaded}
\begin{Highlighting}[]
\CommentTok{\# Examples of different data classes}
\NormalTok{weight\_kg   }\CommentTok{\# numeric, integer, double}
\end{Highlighting}
\end{Shaded}

\begin{verbatim}
[1] 3 2 4 9 7 3 6
\end{verbatim}

\begin{Shaded}
\begin{Highlighting}[]
\NormalTok{animals     }\CommentTok{\# character}
\end{Highlighting}
\end{Shaded}

\begin{verbatim}
[1] "cat" "rat" "bat" "rat"
\end{verbatim}

\begin{Shaded}
\begin{Highlighting}[]
\NormalTok{animal\_size }\OtherTok{\textless{}{-}} \FunctionTok{c}\NormalTok{(}\StringTok{"medium"}\NormalTok{, }\StringTok{"large"}\NormalTok{, }\StringTok{"small"}\NormalTok{, }\StringTok{"medium"}\NormalTok{)}
\NormalTok{animal\_size }\OtherTok{\textless{}{-}} \FunctionTok{factor}\NormalTok{(animal\_size, }\AttributeTok{levels =} \FunctionTok{c}\NormalTok{(}\StringTok{"small"}\NormalTok{, }\StringTok{"medium"}\NormalTok{, }\StringTok{"large"}\NormalTok{))}
\NormalTok{animal\_size  }\CommentTok{\# factor, put in order }
\end{Highlighting}
\end{Shaded}

\begin{verbatim}
[1] medium large  small  medium
Levels: small medium large
\end{verbatim}

\begin{Shaded}
\begin{Highlighting}[]
\NormalTok{logic }\OtherTok{\textless{}{-}} \FunctionTok{c}\NormalTok{(T, F, F, T)  }\CommentTok{\# logical}
\NormalTok{logic}
\end{Highlighting}
\end{Shaded}

\begin{verbatim}
[1]  TRUE FALSE FALSE  TRUE
\end{verbatim}

Vectors have to contain elements that are all of the same class. What
happens if we put data of different classes into one vector?

\begin{Shaded}
\begin{Highlighting}[]
\NormalTok{vec }\OtherTok{\textless{}{-}} \FunctionTok{c}\NormalTok{(}\DecValTok{1}\NormalTok{, }\FloatTok{1.000}\NormalTok{, }\StringTok{"1"}\NormalTok{)}
\CommentTok{\# R converts everything to character because one element is text.}
\end{Highlighting}
\end{Shaded}

\subsection{Subsetting Vectors}\label{subsetting-vectors}

Sometimes we want to keep specific values from a vector. This is called
subsetting (taking a smaller set of the original).

We can subset vectors in two different ways:

\begin{itemize}
\tightlist
\item
  by index (position)
\item
  by condition
\end{itemize}

Regardless of which type of subsetting we choose, we indicate that we
want to subset by using square brackets: \texttt{{[}{]}}.

\subsubsection{Subsetting by Index}\label{subsetting-by-index}

When we subset by index, we are subsetting based on the position of an
element in the vector.

\begin{Shaded}
\begin{Highlighting}[]
\CommentTok{\# Use square brackets}
\NormalTok{weight\_kg[}\DecValTok{2}\NormalTok{]    }\CommentTok{\# returns the 2nd element in the vector}
\end{Highlighting}
\end{Shaded}

\begin{verbatim}
[1] 2
\end{verbatim}

\begin{Shaded}
\begin{Highlighting}[]
\NormalTok{weight\_kg[}\DecValTok{2}\SpecialCharTok{:}\DecValTok{4}\NormalTok{]  }\CommentTok{\# returns the 2nd, 3rd, and 4th elements in the vector}
\end{Highlighting}
\end{Shaded}

\begin{verbatim}
[1] 2 4 9
\end{verbatim}

\subsubsection{Subsetting by Condition}\label{subsetting-by-condition}

Sometimes we don't know or don't want to list out all of the locations
for the data we need. Instead, we might want to subset based on a
quality of the data itself.

To do this, we set a ``condition'' that must be met in order for the
data to be returned.

\begin{Shaded}
\begin{Highlighting}[]
\CommentTok{\# let\textquotesingle{}s start with a condition}
\NormalTok{weight\_kg }\SpecialCharTok{\textgreater{}} \DecValTok{5}
\end{Highlighting}
\end{Shaded}

\begin{verbatim}
[1] FALSE FALSE FALSE  TRUE  TRUE FALSE  TRUE
\end{verbatim}

\begin{Shaded}
\begin{Highlighting}[]
\CommentTok{\# this returns a logical vector (true/false) that tells us which positions meet the condition.}

\CommentTok{\# we now put that condition inside the square brackets}
\NormalTok{weight\_kg[weight\_kg }\SpecialCharTok{\textgreater{}} \DecValTok{5}\NormalTok{]}
\end{Highlighting}
\end{Shaded}

\begin{verbatim}
[1] 9 7 6
\end{verbatim}

\begin{Shaded}
\begin{Highlighting}[]
\CommentTok{\# we can also do this with characters}
\NormalTok{animals }\SpecialCharTok{==} \StringTok{"cat"}
\end{Highlighting}
\end{Shaded}

\begin{verbatim}
[1]  TRUE FALSE FALSE FALSE
\end{verbatim}

\begin{Shaded}
\begin{Highlighting}[]
\NormalTok{animals[animals }\SpecialCharTok{==} \StringTok{"cat"}\NormalTok{]}
\end{Highlighting}
\end{Shaded}

\begin{verbatim}
[1] "cat"
\end{verbatim}

\subsection{Functions}\label{functions}

\ul{Functions} are pre-written bits of codes that perform specific tasks
for us. Functions are always followed by parentheses.

Anything you type into the parentheses are called \ul{arguments}.
Arguments are pieces of information that we give to a function so it
performs its task the way we want it to. To add more than one argument,
you separate them with a comma.

\begin{Shaded}
\begin{Highlighting}[]
\DocumentationTok{\#\# Functions}
\NormalTok{weight\_kg\_mean }\OtherTok{\textless{}{-}} \FunctionTok{mean}\NormalTok{(weight\_kg)   }\CommentTok{\# average of the mass\_kg vector from above}
\NormalTok{weight\_kg\_mean}
\end{Highlighting}
\end{Shaded}

\begin{verbatim}
[1] 4.857143
\end{verbatim}

\begin{Shaded}
\begin{Highlighting}[]
\CommentTok{\# separate arguments with commas}
\FunctionTok{round}\NormalTok{(weight\_kg\_mean)             }\CommentTok{\# rounding}
\end{Highlighting}
\end{Shaded}

\begin{verbatim}
[1] 5
\end{verbatim}

\begin{Shaded}
\begin{Highlighting}[]
\FunctionTok{round}\NormalTok{(weight\_kg\_mean, }\AttributeTok{digits =} \DecValTok{2}\NormalTok{) }\CommentTok{\# round to 2 digits past 0}
\end{Highlighting}
\end{Shaded}

\begin{verbatim}
[1] 4.86
\end{verbatim}

To get more information about a function, use the \texttt{help()}
function or \texttt{?name\_of\_function}.

\begin{Shaded}
\begin{Highlighting}[]
\FunctionTok{help}\NormalTok{(round) }\CommentTok{\# or type ?help}
\end{Highlighting}
\end{Shaded}

\begin{verbatim}
starting httpd help server ... done
\end{verbatim}

We can use a function called \texttt{class()} to figure out the data
type of a vector.

\begin{Shaded}
\begin{Highlighting}[]
\FunctionTok{class}\NormalTok{(weight\_kg)}
\end{Highlighting}
\end{Shaded}

\begin{verbatim}
[1] "numeric"
\end{verbatim}

\subsubsection{Small Group Challenge}\label{small-group-challenge}

Let's practice! Write a few lines of code that do the following:

\begin{itemize}
\tightlist
\item
  create a vector with numbers from 6 to 1 (6, 5, 4, 3, 2, 1)
\item
  assign the vector to an object named \texttt{six\_to\_one}
\item
  subset \texttt{six\_to\_one} to include the last 3 numbers (should
  include 3, 2, 1)
\item
  find the sum of the numbers (hint: use the \texttt{sum()} function)
\end{itemize}

Answer: 6

\begin{Shaded}
\begin{Highlighting}[]
\NormalTok{six\_to\_one }\OtherTok{\textless{}{-}} \FunctionTok{c}\NormalTok{(}\DecValTok{6}\NormalTok{, }\DecValTok{5}\NormalTok{, }\DecValTok{4}\NormalTok{, }\DecValTok{3}\NormalTok{, }\DecValTok{2}\NormalTok{, }\DecValTok{1}\NormalTok{)}
\NormalTok{six\_to\_one}
\end{Highlighting}
\end{Shaded}

\begin{verbatim}
[1] 6 5 4 3 2 1
\end{verbatim}

\begin{Shaded}
\begin{Highlighting}[]
\CommentTok{\# subsetting by index (position)}
\NormalTok{last\_three }\OtherTok{\textless{}{-}}\NormalTok{ six\_to\_one[}\DecValTok{4}\SpecialCharTok{:}\DecValTok{6}\NormalTok{]}
\NormalTok{last\_three}
\end{Highlighting}
\end{Shaded}

\begin{verbatim}
[1] 3 2 1
\end{verbatim}

\begin{Shaded}
\begin{Highlighting}[]
\CommentTok{\# alternate: subsetting by condition}
\NormalTok{last\_three }\OtherTok{\textless{}{-}}\NormalTok{ six\_to\_one[six\_to\_one }\SpecialCharTok{\textless{}} \DecValTok{4}\NormalTok{]}
\NormalTok{last\_three}
\end{Highlighting}
\end{Shaded}

\begin{verbatim}
[1] 3 2 1
\end{verbatim}

\begin{Shaded}
\begin{Highlighting}[]
\FunctionTok{sum}\NormalTok{(last\_three)}
\end{Highlighting}
\end{Shaded}

\begin{verbatim}
[1] 6
\end{verbatim}

Already finished? See if you can condense your code down any further or
turn around and help out a neighbor.

\begin{Shaded}
\begin{Highlighting}[]
\NormalTok{six\_to\_one }\OtherTok{\textless{}{-}} \FunctionTok{seq}\NormalTok{(}\DecValTok{6}\NormalTok{, }\DecValTok{1}\NormalTok{)}
\FunctionTok{sum}\NormalTok{(six\_to\_one[}\DecValTok{4}\SpecialCharTok{:}\DecValTok{6}\NormalTok{])}
\end{Highlighting}
\end{Shaded}

\begin{verbatim}
[1] 6
\end{verbatim}

\subsection{2-Dimensional Data: Data
Frames}\label{dimensional-data-data-frames}

Most of the data you will encounter is two-dimensional (i.e., it has
columns and rows). Its structure resembles a spreadsheet. R is really
good with these types of data. We call these 2D object \ul{data frames}.

\begin{itemize}
\tightlist
\item
  \textbf{rows} go side-to-side
\item
  \textbf{columns} go up-and-down
\end{itemize}

Columns typically represent variables (a factor, trait, or condition) we
are interested in.

Rows represent observations. Each row will be one set of observations.

\pandocbounded{\includegraphics[keepaspectratio]{row_column.png}}

Data frames are made up of multiple vectors. Each vector becomes a
column.

\begin{Shaded}
\begin{Highlighting}[]
\CommentTok{\# Create a simple data frame from scratch}
\NormalTok{plants }\OtherTok{\textless{}{-}} \FunctionTok{data.frame}\NormalTok{(}\AttributeTok{height =} \FunctionTok{c}\NormalTok{(}\DecValTok{55}\NormalTok{, }\DecValTok{17}\NormalTok{, }\DecValTok{42}\NormalTok{, }\DecValTok{47}\NormalTok{, }\DecValTok{68}\NormalTok{, }\DecValTok{39}\NormalTok{, }\DecValTok{51}\NormalTok{, }\DecValTok{23}\NormalTok{),}
                     \AttributeTok{nitrogen =} \FunctionTok{c}\NormalTok{(}\StringTok{"Y"}\NormalTok{, }\StringTok{"N"}\NormalTok{, }\StringTok{"N"}\NormalTok{, }\StringTok{"Y"}\NormalTok{, }\StringTok{"Y"}\NormalTok{, }\StringTok{"N"}\NormalTok{, }\StringTok{"Y"}\NormalTok{, }\StringTok{"N"}\NormalTok{))}

\NormalTok{plants}
\end{Highlighting}
\end{Shaded}

\begin{verbatim}
  height nitrogen
1     55        Y
2     17        N
3     42        N
4     47        Y
5     68        Y
6     39        N
7     51        Y
8     23        N
\end{verbatim}

\subsection{Subsetting Data Frames}\label{subsetting-data-frames}

Because data frames are two-dimensional, we can subset the data in a
data frame by selecting specific columns, specific rows, or both!

R \emph{always} takes information for the row first, then the column.

Think of it as: data{[}row, column{]}

Just like with vectors, we can subset data frames by index or by
condition using square brackets.

The pattern is \texttt{dataframe{[}rows,\ columns{]}}.

\subsubsection{Subsetting by Index}\label{subsetting-by-index-1}

\begin{Shaded}
\begin{Highlighting}[]
\CommentTok{\# Sub{-}setting data frames}
\CommentTok{\# 2{-}dimensional, so you need to specify row and then column}
\CommentTok{\# plants[3] \# doesn\textquotesingle{}t work}

\CommentTok{\# row then column}
\NormalTok{plants[}\DecValTok{4}\NormalTok{,}\DecValTok{1}\NormalTok{]}
\end{Highlighting}
\end{Shaded}

\begin{verbatim}
[1] 47
\end{verbatim}

\begin{Shaded}
\begin{Highlighting}[]
\NormalTok{plants[,}\DecValTok{2}\NormalTok{]}
\end{Highlighting}
\end{Shaded}

\begin{verbatim}
[1] "Y" "N" "N" "Y" "Y" "N" "Y" "N"
\end{verbatim}

Another way to pull out a single column from a data frame is with the
\texttt{\$} operator. This can really come in handy when you know the
name of the column but not the position.

\begin{Shaded}
\begin{Highlighting}[]
\NormalTok{plants}\SpecialCharTok{$}\NormalTok{height}
\end{Highlighting}
\end{Shaded}

\begin{verbatim}
[1] 55 17 42 47 68 39 51 23
\end{verbatim}

\begin{Shaded}
\begin{Highlighting}[]
\CommentTok{\# The $ operator allows you to refer to a column by name instead of position.}
\end{Highlighting}
\end{Shaded}

Regardless of how you specify the column, you can put that code inside
of a function, such as the \texttt{mean()}.

\begin{Shaded}
\begin{Highlighting}[]
\FunctionTok{mean}\NormalTok{(plants}\SpecialCharTok{$}\NormalTok{height)}
\end{Highlighting}
\end{Shaded}

\begin{verbatim}
[1] 42.75
\end{verbatim}

\subsubsection{Subsetting by Condition}\label{subsetting-by-condition-1}

This is a simple data set, but we can use it to ask a question.

Example: Are the heights of plants treated with nitrogen different from
those not treated?

First, we will need to keep only the plants that were treated with
nitrogen.

\begin{Shaded}
\begin{Highlighting}[]
\CommentTok{\# filter rows based on values in the nitrogen column}
\NormalTok{plants[plants}\SpecialCharTok{$}\NormalTok{nitrogen }\SpecialCharTok{==} \StringTok{"Y"}\NormalTok{, ]}
\end{Highlighting}
\end{Shaded}

\begin{verbatim}
  height nitrogen
1     55        Y
4     47        Y
5     68        Y
7     51        Y
\end{verbatim}

\begin{Shaded}
\begin{Highlighting}[]
\CommentTok{\# notice how the comma is still required. leaving it blank after the comma means "keep all columns."}

\CommentTok{\# calculate the mean}
\FunctionTok{mean}\NormalTok{(plants[plants}\SpecialCharTok{$}\NormalTok{nitrogen }\SpecialCharTok{==} \StringTok{"Y"}\NormalTok{, }\DecValTok{1}\NormalTok{])}
\end{Highlighting}
\end{Shaded}

\begin{verbatim}
[1] 55.25
\end{verbatim}

We can create a new data frame by saving the subset data frame to a new
object.

\begin{Shaded}
\begin{Highlighting}[]
\NormalTok{plants\_no\_nitrogen }\OtherTok{\textless{}{-}}\NormalTok{ plants[plants}\SpecialCharTok{$}\NormalTok{nitrogen }\SpecialCharTok{==} \StringTok{"N"}\NormalTok{, ]}
\end{Highlighting}
\end{Shaded}

\subsubsection{Small Group Challenge (5
min)}\label{small-group-challenge-5-min}

As a group, find the standard deviation (\texttt{sd()}) of the height of
plants treated with nitrogen and those not treated with nitrogen. Which
group has the larger standard deviation? Any ideas what that means?

\begin{Shaded}
\begin{Highlighting}[]
\FunctionTok{sd}\NormalTok{(plants[plants}\SpecialCharTok{$}\NormalTok{nitrogen }\SpecialCharTok{==} \StringTok{"Y"}\NormalTok{, }\DecValTok{1}\NormalTok{])}
\end{Highlighting}
\end{Shaded}

\begin{verbatim}
[1] 9.105859
\end{verbatim}

\begin{Shaded}
\begin{Highlighting}[]
\FunctionTok{sd}\NormalTok{(plants[plants}\SpecialCharTok{$}\NormalTok{nitrogen }\SpecialCharTok{==} \StringTok{"N"}\NormalTok{, }\DecValTok{1}\NormalTok{])}
\end{Highlighting}
\end{Shaded}

\begin{verbatim}
[1] 12.14839
\end{verbatim}

\subsection{Helpful Functions}\label{helpful-functions}

Below are some functions that I often find very helpful when working
with vectors and data frames:

\begin{itemize}
\tightlist
\item
  \texttt{str()}: shows the structure of the object (e.g., rows and
  columns)
\item
  \texttt{head()} and \texttt{tail()}: shows the first six and last six
  rows, respectfully
\item
  \texttt{length()}: counts the number of elements in an object
\item
  \texttt{ncol()} and \texttt{nrow()}: counts the number of columns or
  rows, respectfully
\item
  \texttt{names()}: shows the names of the columns in a data frame
\item
  \texttt{unique()}: shows one of each element in an object (removes
  duplicate values)
\end{itemize}

\begin{Shaded}
\begin{Highlighting}[]
\FunctionTok{str}\NormalTok{(plants) }\CommentTok{\# structure of the object}
\end{Highlighting}
\end{Shaded}

\begin{verbatim}
'data.frame':   8 obs. of  2 variables:
 $ height  : num  55 17 42 47 68 39 51 23
 $ nitrogen: chr  "Y" "N" "N" "Y" ...
\end{verbatim}

\begin{Shaded}
\begin{Highlighting}[]
\FunctionTok{head}\NormalTok{(plants) }\CommentTok{\# first 6 values or rows}
\end{Highlighting}
\end{Shaded}

\begin{verbatim}
  height nitrogen
1     55        Y
2     17        N
3     42        N
4     47        Y
5     68        Y
6     39        N
\end{verbatim}

\begin{Shaded}
\begin{Highlighting}[]
\FunctionTok{head}\NormalTok{(plants, }\AttributeTok{n =} \DecValTok{4}\NormalTok{) }\CommentTok{\# first n values or rows}
\end{Highlighting}
\end{Shaded}

\begin{verbatim}
  height nitrogen
1     55        Y
2     17        N
3     42        N
4     47        Y
\end{verbatim}

\begin{Shaded}
\begin{Highlighting}[]
\FunctionTok{tail}\NormalTok{(plants, }\AttributeTok{n =} \DecValTok{4}\NormalTok{) }\CommentTok{\# last n values or rows}
\end{Highlighting}
\end{Shaded}

\begin{verbatim}
  height nitrogen
5     68        Y
6     39        N
7     51        Y
8     23        N
\end{verbatim}

\begin{Shaded}
\begin{Highlighting}[]
\FunctionTok{length}\NormalTok{(plants)  }\CommentTok{\# for a dataframe, number of columns}
\end{Highlighting}
\end{Shaded}

\begin{verbatim}
[1] 2
\end{verbatim}

\begin{Shaded}
\begin{Highlighting}[]
\FunctionTok{length}\NormalTok{(plants}\SpecialCharTok{$}\NormalTok{height) }\CommentTok{\# for a column, number of rows}
\end{Highlighting}
\end{Shaded}

\begin{verbatim}
[1] 8
\end{verbatim}

\begin{Shaded}
\begin{Highlighting}[]
\FunctionTok{ncol}\NormalTok{(plants)  }\CommentTok{\# number of columns}
\end{Highlighting}
\end{Shaded}

\begin{verbatim}
[1] 2
\end{verbatim}

\begin{Shaded}
\begin{Highlighting}[]
\FunctionTok{nrow}\NormalTok{(plants)  }\CommentTok{\# number of rows}
\end{Highlighting}
\end{Shaded}

\begin{verbatim}
[1] 8
\end{verbatim}

\begin{Shaded}
\begin{Highlighting}[]
\FunctionTok{names}\NormalTok{(plants) }\CommentTok{\# list of column or object names}
\end{Highlighting}
\end{Shaded}

\begin{verbatim}
[1] "height"   "nitrogen"
\end{verbatim}

\begin{Shaded}
\begin{Highlighting}[]
\FunctionTok{unique}\NormalTok{(plants}\SpecialCharTok{$}\NormalTok{nitrogen) }\CommentTok{\# one of each value present in the column, duplicates removed}
\end{Highlighting}
\end{Shaded}

\begin{verbatim}
[1] "Y" "N"
\end{verbatim}




\end{document}
